% !TEX root = ../../../main.tex
% 来源：中学数学实验教材·第二册（几何）
% 章节：圆 / 圆与直线的位置关系
% 原始文件：4.tex，行 1767-1793
% 类型：tikz
% ---
    \begin{tikzpicture}[>=latex, scale=.7]
  \tkzDefPoints{0/0/I}\tkzDrawPoints(I)
  \tkzDefPoint(15:1.5){E}
  \tkzDefPoint(135:1.5){F}
  \tkzDefPoint(-90:1.5){D}
  \tkzLabelPoints[right](I)
  \tkzDefTangent[at=E](I) \tkzGetPoint{a'}
  \tkzDefTangent[at=F](I) \tkzGetPoint{b'}
  \tkzDefTangent[at=D](I) \tkzGetPoint{c'}
  \tkzInterLL(E,a')(F,b')\tkzGetPoint{A}
  \tkzInterLL(D,c')(F,b')\tkzGetPoint{B}
  \tkzInterLL(D,c')(E,a')\tkzGetPoint{C}
  \tkzLabelPoints[below](B,D,C)
  \tkzLabelPoints[above](A)
  \tkzLabelPoints[left](F)
  \tkzLabelPoints[right](E)
  \tkzDrawPolygon[thick](A,B,C)
  \tkzDrawCircle[thick](I,E)
  \tkzLabelSegment[below](B,D){$y$}
  \tkzLabelSegment[below](C,D){$z$}
  \tkzLabelSegment[left](B,F){$y$}
  \tkzLabelSegment[left](A,F){$x$}
  \tkzLabelSegment[right](A,E){$x$}
  \tkzLabelSegment[right](E,C){$z$}
  
  
    \end{tikzpicture}
